\begin{helper}[Uniqueness of the enlarged-set crossing]
In the setting of \noderef{ScheideLastExitPath}, fix a defective index
$\delta$ and suppose every vertex-node of $L.path$ through
$L.\mathit{lastInside}$ belongs to the replacement-tree vertex set
$V((R\delta).\mathit{replacementTree})$.  Then every edge identity that
belongs both to
\noderef{BoundaryEdges}$(G,\noderef{ScheideEnlargedVertexSet}(R))$ and to
the full $(D.\alpha_0,\alpha)$ component through $y_0$ equals
$L.\mathit{crossing}$.
\end{helper}
\begin{proof}
Put $X=\noderef{ScheideEnlargedVertexSet}(R)$ and
$v=L.\mathit{lastInside}$.  The exact replacement vertex-set field in
\noderef{ScheideReplacementData} gives
$V((R\delta).\mathit{replacementTree})\subseteq X$.  The assumed prefix
confinement therefore puts every vertex-node of $L.path$ through $v$ in
$X$.  Conversely, the $afterLastOutside$ field of
\noderef{ScheideLastExitPath} puts every later vertex-node of the path
outside $X$.

Let $f$ be both a boundary edge of $X$ and an edge-node of the indicated
two-color component.  The support equality $L.pathSpans$ from
\noderef{ScheideLastExitPath}, expanded through
\noderef{TwoColorComponent}, places $\operatorname{inr}(f)$ on $L.path$.
The walk alternates vertex-nodes and edge-nodes because adjacency in
\noderef{TwoColorIncidenceGraph} only joins the two summands.  Thus the
two path neighbors of this occurrence of $f$ are precisely its two
endpoints in $G$.  Since $f\in\noderef{BoundaryEdges}(G,X)$, exactly one
of these endpoint vertices lies in $X$.

If the occurrence of $f$ were before the edge immediately following
$v$, both neighboring vertex-nodes would occur no later than $v$ and
hence both would lie in $X$.  If it were after that edge, both would
occur after $v$ and hence both would lie outside $X$.  Both alternatives
contradict that $f$ is a boundary edge.  The only remaining occurrence is
the edge-node immediately following $v$.  The path-simplicity field of
\noderef{ScheideLastExitPath} makes this occurrence unique, and its
$crossingAt$ field identifies it as $L.\mathit{crossing}$.  Therefore
$f=L.\mathit{crossing}$.
\end{proof}
